\documentclass[12pt]{article}
\usepackage[T1]{fontenc}
\usepackage[utf8]{inputenc}
\usepackage{lmodern}
\usepackage{textcomp}
\usepackage{enumitem}
\usepackage{geometry}
\geometry{margin=1in}
\begin{document}
\begin{center}
Answer Key - History - Undergraduate Level Assessment 1 \
\small (Answers are ordered exactly as in the question paper.)
\end{center}

\section*{Multiple Choice}
\begin{enumerate}[label=\arabic*.]
\item \textbf{Answer:} B
\\ \textit{Options:} A) 5,000,000; 100\%; B) 3,000,000; 75\%; C) 4,000,000; 80\%; D) 2,500,000; 90\%; E) 1,500,000; 70\%
\\ \textit{Note:} correct because historical estimates suggest that ancient Egypt had a population of around 3,000,000 people, with approximately 75\% of them engaged in agriculture, highlighting the central role of farming in the economy and society of the time.
\item \textbf{Answer:} A
\\ \textit{Options:} A) North America; B) Australia; C) Central America; D) Arctic Region; E) Middle East
\\ \textit{Note:} The Hoabinhian tradition, characterized by the production of stone tools from chipped pebbles, is primarily associated with prehistoric cultures in Southeast Asia, not North America; thus, the correct answer should reflect the geographical origins of this archaeological phenomenon.
\item \textbf{Answer:} D
\\ \textit{Options:} A) trade.; B) warfare.; C) the development of a common religion.; D) conspicuous consumption by elites.
\\ \textit{Note:} The Minoan culture, characterized by its advanced trade and artistic achievements, exhibits limited evidence of conspicuous consumption by elites, suggesting a more egalitarian society compared to other early civilizations that displayed clear social hierarchies and extravagant displays of wealth.
\item \textbf{Answer:} E
\\ \textit{Options:} A) music and arts.; B) geography and climate.; C) architecture and construction methods.; D) legal and judicial systems.; E) political and military history.
\\ \textit{Note:} The correct answer is supported by the fact that Maya records, such as stelae and codices, predominantly document the reigns of rulers, military conquests, and political events, whereas Mesopotamian writings include a broader range of topics, including economic transactions and religious texts.
\item \textbf{Answer:} A
\\ \textit{Options:} A) external; extrafamilial; B) interfamilial; extrafamilial; C) religious; military; D) internal; intrafamilial; E) external; internal
\\ \textit{Note:} The correct answer highlights that conflict models focus on external tensions arising from competition and power struggles between societies, while integration models prioritize extrafamilial tensions that facilitate social cohesion and cooperation within societies.
\end{enumerate}

\section*{True / False}
\begin{enumerate}[label=\arabic*.]
\setcounter{enumi}{5}
\item \textbf{Answer:} True
\\ \textit{Options:} A) True; B) False
\\ \textit{Note:} According to the passage: 20\%
\item \textbf{Answer:} False
\\ \textit{Options:} A) True; B) False
\\ \textit{Note:} The correct answer is: Interstate 95
\item \textbf{Answer:} True
\\ \textit{Options:} A) True; B) False
\\ \textit{Note:} According to the passage: Spiracles
\item \textbf{Answer:} False
\\ \textit{Options:} A) True; B) False
\\ \textit{Note:} The correct answer is: regional genres of the Venezuelan Andes
\item \textbf{Answer:} False
\\ \textit{Options:} A) True; B) False
\\ \textit{Note:} The correct answer is: accession into NATO and the European Union
\end{enumerate}

\section*{Long Form Response}
\begin{enumerate}[label=\arabic*.]
\setcounter{enumi}{10}
\item \textbf{Answer:} Orthodox church
\\ \textit{Note:} Expected answer: Orthodox church
\item \textbf{Answer:} 1400
\\ \textit{Note:} Expected answer: 1400
\end{enumerate}

\vfill
\noindent\textit{End of Answer Key}
\end{document}